% !TeX program = xelatex
\documentclass[12pt,a4paper]{article}
\usepackage[margin=2cm]{geometry}
\usepackage{fontspec}
\usepackage{amsmath,amssymb}
\usepackage{enumitem}
\setmainfont{Times New Roman}
\setlength{\parindent}{0pt}
\setlength{\parskip}{5pt}
\begin{document}

\begin{center}
{\large\bfseries SỞ GIÁO DỤC VÀ ĐÀO TẠO HÀ NỘI}\par
{\Large\bfseries KỲ THI TUYỂN SINH LỚP 10 THPT}\par
{\large NĂM HỌC 2021--2022}\par
\vspace{4pt}
{\bfseries MÔN THI: TOÁN}\par
Ngày thi: 13/6/2021 \qquad Thời gian làm bài: 90 phút
\end{center}

\textit{Phần trích theo phạm vi Đại số, toán thực tế và diện tích/thể tích. Điểm từng câu giữ theo đề gốc; tổng điểm phần trích là 7,0 điểm trên thang điểm gốc.}

\textbf{Bài I (2,0 điểm)}

Cho hai biểu thức
\[
A=\frac{\sqrt{x}}{\sqrt{x}+3},\qquad
B=\frac{2\sqrt{x}}{\sqrt{x}-3}-\frac{3x+9}{x-9},
\]
với $x\geq 0$, $x\neq 9$.

\begin{enumerate}[label=\arabic*)]
  \item Tính giá trị của biểu thức $A$ khi $x=16$.
  \item Chứng minh rằng
  \[
  A+B=\frac{3}{\sqrt{x}+3}.
  \]
\end{enumerate}

\textbf{Bài II (2,5 điểm)}

\begin{enumerate}[label=\arabic*)]
  \item Giải bài toán sau bằng cách lập phương trình hoặc hệ phương trình:

  Một tổ sản xuất phải làm xong 4800 bộ đồ bảo hộ y tế trong một số ngày quy định. Thực tế, mỗi ngày tổ đó đã làm được nhiều hơn 100 bộ đồ bảo hộ y tế so với số bộ đồ bảo hộ y tế phải làm trong một ngày theo kế hoạch. Vì thế 8 ngày trước khi hết thời hạn, tổ sản xuất đã làm xong 4800 bộ đồ bảo hộ y tế đó. Hỏi theo kế hoạch, mỗi ngày tổ sản xuất phải làm bao nhiêu bộ đồ bảo hộ y tế? Giả định rằng số bộ đồ bảo hộ y tế mà tổ đó làm xong trong mỗi ngày là bằng nhau.

  \item Một thùng nước có dạng hình trụ với chiều cao $1{,}6\,\mathrm{m}$ và bán kính đáy $0{,}5\,\mathrm{m}$. Người ta sơn toàn bộ phía ngoài mặt xung quanh của thùng nước này (trừ hai mặt đáy). Tính diện tích bề mặt được sơn của thùng nước (lấy $\pi\approx 3{,}14$).
\end{enumerate}

\textbf{Bài III (2,0 điểm)}

\begin{enumerate}[label=\arabic*)]
  \item Giải hệ phương trình
  \[
  \begin{cases}
  \dfrac{3}{x+1}-2y=-1,\\[4pt]
  \dfrac{5}{x+1}+3y=11.
  \end{cases}
  \]
  \item Trong mặt phẳng tọa độ $Oxy$, cho parabol $(P): y=x^2$ và đường thẳng $(d): y=2x+m-2$. Tìm tất cả giá trị của $m$ để $(d)$ cắt $(P)$ tại hai điểm phân biệt có hoành độ $x_1,x_2$ sao cho $|x_1-x_2|=2$.
\end{enumerate}

\textbf{Bài V (0,5 điểm)}

Với các số thực $a$ và $b$ thỏa mãn $a^2+b^2=2$, tìm giá trị nhỏ nhất của biểu thức
\[
P=3(a+b)+ab.
\]

\end{document}
